\section{Exact periodic scheduling with quadratic merger work}
\label{sec:queue}

\subsection{Interval pairings and the unchanged merger rule}

An interval pairing identifies the integer points of two intervals of equal
length by an isometry. We use inclusive endpoints in this section:
\[
 g:[a,b]\longrightarrow[c,d],\qquad a\le c,\quad b-a=d-c.
\]
The map is either \(x\mapsto x+c-a\) or the reflection
\(x\mapsto a+d-x\). A periodic pairing is an increasing pairing with
\(a<c\le b+1\). Its period is \(p=c-a>0\), and its support is \([a,d]\).
The convention includes two adjacent intervals.

For periodic pairings with periods \(p,q\), write \(h=\gcd(p,q)\).
The maintained sharp rule merges them when their support overlap has at least
\[
 p+q-h
\]
integer points. Their union is then represented by a period-\(h\)
pairing. This is the Fine--Wilf periodicity criterion in the form already
implemented and verified by the repository \cite{FineWilf1965,ProveIt2026}.
The alternative classical threshold \(p+q\) remains available. Our
scheduler uses the existing arithmetic test unchanged. The arguments below
about scheduling work for either rule.

The historical algorithm scans pairs of list indices in lexicographic order.
It accepts the first legal pair \((i,j)\), keeps the replacement at \(i\),
deletes \(j\), and restarts from the beginning. Since a merge changes only
one surviving row, most failed tests in the repeated prefix compare
unchanged data.

\subsection{Stable identifiers and generations}

Give each input row a permanent identifier equal to its original list
position. Active identifiers remain in increasing order when rows are
deleted. For each identifier keep a generation counter, initially zero.
Whenever a row is replaced, increment its counter. A heap entry contains
\[
 (i,j,\nu_i,\nu_j,\text{replacement}),
\]
ordered first by \(i\), then by \(j\). Its generations record the operands
on which the exact merger test was performed.

Initialisation tests every pair of periodic rows and inserts the successful
pairs. To perform a merge, discard stale heap minima until the first
current entry is found. An entry is stale when either identifier has
disappeared or either generation has changed. Replace the lower-identifier
row, remove the other, and test exactly the pairs incident to the
replacement. Entries for untouched pairs remain valid. No equality test
uses a fingerprint or a probabilistic hash.

\begin{theorem}[Exact simulation of the historical closure]
\label{thm:queue-exact}
For every finite ordered pairing list and either maintained periodic
threshold, the queue closure performs exactly the historical sequence of
successful merges. It returns the same ordered list. When it records current
list indices, its merger certificate is identical to the historical
certificate.
\end{theorem}
\begin{proof}
Maintain the following invariant: each legal current pair has exactly one
nonstale heap entry, and every nonstale entry is the correct successful test
on its present operands. Initialisation establishes the invariant.

Suppose it holds before a merge of \(i<j\). Pairs avoiding \(i,j\) have
unchanged operands, so both their legality and their existing entries are
unchanged. Deletion of \(j\) invalidates every entry involving \(j\).
Incrementing the generation of \(i\) invalidates all old entries involving
\(i\). Testing each surviving partner of the new \(i\) restores precisely
the missing current entries. The invariant is preserved.

The stable identifiers have the same relative order as current list
positions. The first nonstale heap entry is therefore the
lexicographically first legal pair of current rows, which is exactly the
historical choice. Induction proves equality of successful choices and
replacements. The rank of an identifier among active identifiers is its
current list index, giving the same certificate indices.
\end{proof}

The producer currently computes these ranks by scanning the active
identifiers. There are at most \(k-1\) mergers, so this adds \(O(k^2)\)
operations over a closure. A separate order-statistics structure would
complicate the implementation without improving its displayed asymptotic
bound.

\begin{corollary}[Complete trace preservation]
\label{cor:trace}
Replacing only periodic-closure scheduling in the maintained AHT routine
preserves the orbit count, the number of AHT cycles, and every complete
local-reduction certificate.
\end{corollary}
\begin{proof}
After each periodic closure the ordered pairing list is identical by
\cref{thm:queue-exact}. All subsequent carrier selection, maximal-power
transmission, suffix truncation, and next-cycle preparation are unchanged
deterministic operations on that list. Induct over the cycles.
\end{proof}

Callback invocation counts and scheduler work statistics are intentionally
not certificate contents. A wall-clock deadline can stop the algorithms at
different points because their costs differ. The corollary concerns
completed mathematical traces; a shared cycle cap still counts the same
begun AHT cycles.

\subsection{Work bound and the adaptive default}

\begin{theorem}[Quadratic candidate count]
\label{thm:queue-work}
For a closure starting with \(k\ge1\) rows, pure queue mode performs at most
\((k-1)^2\) exact merger tests. It requires at most \(O(k^2)\) heap entries.
For endpoint bit length \(B\), a conservative bit bound is
\[
 O\!\left(k^2\bigl(G(B)+B+\log^2(k+1)\bigr)\right),
\]
where \(G(B)\) bounds a gcd and the bounded number of endpoint operations
used in one merger test. The empty list requires no test.
\end{theorem}
\begin{proof}
Initialisation uses at most \(\binom{k}{2}\) tests. A merge leaving \(r\)
rows creates at most \(r-1\) incident candidates. Since the live row count
strictly decreases,
\[
 \binom{k}{2}+\sum_{r=1}^{k-1}(r-1)
 =\binom{k}{2}+\binom{k-1}{2}
 =(k-1)^2.
\]
Nonperiodic rows only reduce the number of candidates. Every successful
test inserts one entry; every entry can be popped at most once, including
stale entries. Thus the total heap size, insertions, and removals are
quadratic. Each heap operation takes \(O(\log(k+1))\) identifier
comparisons, on \(O(\log(k+1))\)-bit keys. Endpoint arithmetic and stored
replacement construction contribute the other terms. Rank scans and
output construction fit within the same bound.
\end{proof}

The storage bound in bits is \(O(k^2(B+\log(k+1)))\), excluding the input
and any separately retained full trace. This is a real tradeoff:
historical restarting needs only a linear working list.

Moreover, an eager queue is unnecessary on easy inputs. If every row is
identical, the historical closure succeeds immediately on its first pair
\(k-1\) times, whereas the queue begins by checking all
\(\binom{k}{2}\) pairs.

The adaptive default therefore first follows the historical candidate
sequence. At each restart it discovers the first two periodic indices and
tests their pair before collecting further indices. After a failure it
finishes indexing the periodic rows once, avoiding repeated inner loops
over nonperiodic rows. It allows at most
\(\binom{k}{2}\) merger tests in this mode. Only an additional pending test
switches to the queue on the current residual list.

\begin{corollary}[Adaptive bound]
\label{cor:adaptive}
Adaptive closure performs at most
\[
 \binom{k}{2}+\max(k-1,0)^2
\]
merger tests, preserves the historical trace, and has the same quadratic
worst-case storage bound as the queue. With no switch, its extra index
storage is linear. No-merge closures and successive first-pair successes
avoid heap construction.
\end{corollary}
\begin{proof}
The preliminary allowance gives the first term. At the switch, at most
\(k\) rows remain, so \cref{thm:queue-work} gives the second.
Both parts perform the same greedy choices. At most \(k\) index
reconstructions, each scanning at most \(k\) rows, add \(O(k^2)\) work.
When no merge succeeds, one pass uses at most the allowance. When the
first tested pair always succeeds, there are only \(k-1\) tests.
\end{proof}

This allowance is an internal scheduling threshold, not an incompleteness
budget. It changes how the answer is found, not whether computation may
stop with a mathematical verdict.

\subsection{A complete five-cycle family}

A per-phase improvement could be hidden by equally expensive later
phases. The following family rules out that explanation for this
particular bottleneck.

\begin{theorem}[An exact cubic-to-quadratic complete-run witness]
\label{thm:five-cycle}
Let \(m\ge2\), \(P=10m\), and \(N=23m\). Define
\[
 g_i:[i,P+2i-1]\longrightarrow[P+2i,2P+3i-1],
 \qquad 0\le i\le m.
\]
Use the ordered list \(g_0,\ldots,g_{m-1}\), followed by \(m\) copies
of \(g_m\). Under the sharp periodic rule, the maintained algorithm
uses exactly five AHT cycles and returns one orbit.
Its historical merger-test count is
\[
 T_{\rm old}(m)=m^3+\frac{m^2+5m}{2}-4,
\]
whereas pure queue mode uses
\[
 T_{\rm queue}(m)=5m^2-6m+2.
\]
Adaptive mode uses \(O(m^2)\) tests. The entire revised count takes
\(O(m^2\polylog m)\) bit operations on this family.
\end{theorem}
\begin{proof}
The period of \(g_i\) is \(P+i\). For \(i<j\), their overlap contains
\(2P+3i-j\) points. A sharp merge would require
\[
 2P+3i-j\ge2P+i+j-\gcd(P+i,P+j),
\]
or \(\gcd(P+i,P+j)\ge2(j-i)\). The gcd divides \(j-i\), so this is
impossible. Identical tail copies merge. The supports cover the universe,
so initially there is no static contraction.

\emph{Cycle 1.}
With \(r\) tail copies remaining, all \(\binom m2\) pairs among the first
\(m\) rows and all \(mr\) mixed pairs fail. If \(r\ge2\), the next tail
pair succeeds. Including the last entirely failed pass, the historical
closure takes
\[
 \sum_{r=1}^{m}\left(\binom m2+mr\right)+(m-1)
 =m^3+m-1
\]
tests. The queue takes \((7m^2-7m+2)/2\).
The carrier is \(g_m\); no other range is contained in its range.
Truncating three final points leaves
\[
 h_m:[m,P+2m-4]\longrightarrow[P+2m,2P+3m-4]
\]
and universe size \(N-3\). This pairing is nonperiodic.

\emph{Cycle 2.}
The \(m\) untouched original pairings remain pairwise unmergeable.
Both modes use \(\binom m2\) tests. The carrier
\(g_{m-1}\) ties \(h_m\) at the right endpoint but has the earlier range
start. One inverse transmission maps \(h_m\) to
\[
 h:[m,P+2m-4]\longrightarrow[m+1,P+2m-3],
\]
which has period one. No other original range moves. Truncating another
three points changes \(g_{m-1}\) to the nonperiodic pairing
\[
 q:[m-1,P+2m-6]\longrightarrow[P+2m-2,2P+3m-7].
\]
The universe now has size \(N-6\).

\emph{Cycle 3.}
The first legal merge is \(g_0\) with \(h\): their overlap has
\(P+m-2\ge P\) points, meeting its sharp threshold. The replacement
has period one and support \([0,2P-1]\). It immediately merges with
\(g_1\), then with \(g_2\), through \(g_{m-2}\). Indeed, the
period-one hull through \(g_j\) overlaps \(g_{j+1}\) in
\(2P+2j-1\) points, at least \(P+j+1\). The final carrier is
\[
 u:[0,N-8]\longrightarrow[1,N-7].
\]
The old closure uses \(m-1\) tests before the first success and \(m-2\)
further immediate successes, or \(2m-3\) altogether. The queue starts
with \(m\) periodic rows and merges all of them, using \((m-1)^2\) tests.

Both the domain and range of \(q\) lie in the range of \(u\).
Their maximal inverse powers are \(m-1\) and \(P+2m-2\),
respectively, so transmission turns \(q\) into the identity on
\([0,P+m-5]\). Truncation leaves universe size \(P+m-4\).

\emph{Cycles 4 and 5.}
The fourth cycle deletes this identity and uses the period-one carrier
to leave one point. The fifth contracts that isolated point, yielding
one orbit. Neither cycle performs a merger test.

Adding the three displayed phase counts gives the two exact formulas.
Every other operation occurs in five cycles on at most \(2m\) rows.
The coordinate bit length is \(O(\log m)\), since \(N=23m\).
The queue and adaptive work bounds therefore give
\(O(m^2\polylog m)\) bit operations for the full count.
\end{proof}

This is a lower bound for the historical \emph{implementation}, not for
orbit counting. Nor is the family claimed to arise from normal surfaces
in knot exteriors. Its role is precise: it demonstrates that a real
redundancy in a maintained component kernel can dominate a complete
run and that the new scheduler removes it.

\subsection{What remains unchanged in the global AHT analysis}

For general systems, the number of cycles is governed by the classical
AHT argument. If a run has \(Q\) cycles, the new closure work is bounded by
\[
 O\!\left(Q k^2\bigl(G(B)+B+\log^2(k+1)\bigr)\right).
\]
One may insert a classical polynomial bound for \(Q\) to obtain a coarse
explicit polynomial estimate. We do not optimize that global exponent
or advertise this as a faster abstract AHT theorem.

The exact simulation theorem is the important compatibility property:
whatever geometric information or independently checkable local trace the
existing code produces is retained. In particular, a future weighted
replayer can consume the same trace. This is the bridge to the next two
sections.
